% Seccion: Introducción general
\section{Introducción general}
\label{sec:cap01_introduccion}

Este apartado expone detalladamente los aspectos fundamentales de introducción general para el colector solar de aire.
El análisis del comportamiento termo-fluidodinámico requiere la formulación de ecuaciones constitutivas específicas.
Se consideran condiciones de frontera realistas para aproximar el modelo a un escenario industrial real.
Los parámetros resultantes serán fundamentales para el proceso de validación cruzada entre modelos teóricos y numéricos.

\begin{cajainfo}{Contexto de Ingeniería}
El diseño de colectores solares requiere un equilibrio riguroso entre la eficiencia de transferencia de calor y las pérdidas de carga en el conducto.
El uso de serpentines busca maximizar la turbulencia y la superficie expuesta.
\end{cajainfo}
